\chapter{Migrating to Bijou 7.1.0}\label{record-271}

\noindent{\footnotesize\color{BookMuted}\path{docs/releases/7.1.0/migration-guide.md}\par}\smallskip

Bijou 7.1.0 is a minor release from 7.0.0. It is intended to be additive for applications using documented public APIs.

\section{Recommended Upgrade}\label{record-271-recommended-upgrade}

Update the runtime packages your app uses together:

\begin{CodeBlock}
npm install \
  @flyingrobots/bijou@7.1.0 \
  @flyingrobots/bijou-node@7.1.0 \
  @flyingrobots/bijou-tui@7.1.0 \
  @flyingrobots/bijou-tui-app@7.1.0
\end{CodeBlock}

If your app uses localization tooling, keep those packages in lock-step too:

\begin{CodeBlock}
npm install \
  @flyingrobots/bijou-i18n@7.1.0 \
  @flyingrobots/bijou-i18n-tools@7.1.0 \
  @flyingrobots/bijou-i18n-tools-node@7.1.0 \
  @flyingrobots/bijou-i18n-tools-xlsx@7.1.0
\end{CodeBlock}

If your app uses the Bijou MCP server, keep it lock-step too:

\begin{CodeBlock}
npm install @flyingrobots/bijou-mcp@7.1.0
\end{CodeBlock}

\section{Runtime Behavior}\label{record-271-runtime-behavior}

No breaking runtime migration is expected.

Applications may notice better behavior in these areas:

\begin{itemize}
\tightlist
\item
  full-screen TUIs repaint correctly after terminal resize, laptop sleep/wake, or external-display geometry changes
\item
  ANSI parsing correctly handles scoped foreground/background resets
\item
  DOGFOOD Blocks guide navigation can focus parent preview rows
\item
  quit confirmation copy is less redundant
\end{itemize}

If your app had local workarounds for stale cells after resize or manual ANSI reset cleanup, retest those paths against 7.1.0 and remove the workaround if Bijou now owns the behavior.

\section{New Proof APIs}\label{record-271-new-proof-apis}

The new \texttt{u\allowbreak{}i\allowbreak{}-\allowbreak{}s\allowbreak{}c\allowbreak{}e\allowbreak{}n\allowbreak{}e\allowbreak{}-\allowbreak{}i\allowbreak{}r\allowbreak{}/\allowbreak{}1} and GraphQL block helpers are available for early proof work, but they are still the first slices of a larger Runtime Graph direction. Use them when you want deterministic receipts, source maps, lower modes, and agent-readable debug facts. Do not assume they are the final V8 product contract yet.

\section{Theme Work}\label{record-271-theme-work}

Theme-token APIs are more useful in 7.1.0, including first-party dark/light presets and safe-pair declarations. Existing raw-color code should keep working. New code should prefer semantic token refs when the color is part of a themeable product surface.

\section{Release Tooling}\label{record-271-release-tooling}

Maintainers preparing releases should use:

\begin{CodeBlock}
npm run release:readiness -- --milestone v7.1.0
\end{CodeBlock}

The milestone-aware mode blocks missing release evidence packets, open target issues, lingering \texttt{w\allowbreak{}o\allowbreak{}r\allowbreak{}k\allowbreak{}-\allowbreak{}i\allowbreak{}n\allowbreak{}-\allowbreak{}p\allowbreak{}r\allowbreak{}o\allowbreak{}g\allowbreak{}r\allowbreak{}e\allowbreak{}s\allowbreak{}s} labels, stale release posture docs, and missing package-smoke coverage before the local gauntlet runs.

DOGFOOD contributors should also expect \texttt{n\allowbreak{}p\allowbreak{}m\allowbreak{} \allowbreak{}r\allowbreak{}u\allowbreak{}n\allowbreak{} \allowbreak{}d\allowbreak{}o\allowbreak{}g\allowbreak{}f\allowbreak{}o\allowbreak{}o\allowbreak{}d\allowbreak{}:\allowbreak{}i\allowbreak{}1\allowbreak{}8\allowbreak{}n\allowbreak{}:\allowbreak{}d\allowbreak{}e\allowbreak{}b\allowbreak{}t} to scan new \texttt{e\allowbreak{}x\allowbreak{}a\allowbreak{}m\allowbreak{}p\allowbreak{}l\allowbreak{}e\allowbreak{}s\allowbreak{}/\allowbreak{}d\allowbreak{}o\allowbreak{}c\allowbreak{}s\allowbreak{}/\allowbreak{}*\allowbreak{}*\allowbreak{}/\allowbreak{}*\allowbreak{}.\allowbreak{}t\allowbreak{}s} modules by default. Tooling-only exclusions must be explicit.

\section{Compatibility Stance}\label{record-271-compatibility-stance}

Bijou 7.1.0 should be a straightforward upgrade from 7.0.0 for application code. The main migration work is optional cleanup: remove local resize/ANSI workarounds, adopt theme tokens where they clarify product intent, and use the new proof helpers only where deterministic UI evidence is valuable.

